\documentclass[11pt]{amsart}
\usepackage{amsmath,amssymb}
\usepackage{xcolor}
\usepackage[margin=1.15in]{geometry}
\usepackage[colorlinks=true,linkcolor=blue!50!black,citecolor=blue!50!black,urlcolor=blue!50!black]{hyperref}
\title{Overview of the collection}
\author{Claude Opus 5.5, GPT-6-Astra, Claude Fable 5.1}
\date{October 29, 2026}
\begin{document}
\begin{abstract}
This collection contains three manuscripts, each resolving a published open question.
All main results are formally verified in Lean~4 with Mathlib.
The manuscripts, Lean code and supporting computations were produced by AI models working as a team of agents; see the repository README for how the results were produced.
\end{abstract}
\maketitle

\section*{001. The maximum number of non-equivalent ancestral configurations}
\emph{Discipline: mathematical biology (phylogenetics); enumerative combinatorics.}

In the multispecies coalescent, gene tree probabilities are computed by dynamic programming over \emph{ancestral configurations}.
Disanto and Rosenberg (Bull.\ Math.\ Biol.\ 2019) introduced the number $c^*(T)$ of non-equivalent root configurations of a matching gene tree / species tree pair of shape $T$ and conjectured that
\[
\max_{|T|=n}c^*(T)=\Bigl\lfloor 3^{\lfloor n/3\rfloor}\,\tfrac{2r^2+8r+23}{27}\Bigr\rfloor,\qquad r=n \bmod 3,\ n\ge 2,
\]
having verified it for $n\le 20$.
The manuscript proves the conjecture for all $n$, shows that the maximizing shape is unique for $n\ge 9$ (it is obtained by repeatedly joining a three-leaf tree), determines the growth constant $3^{1/3}$ (previously known to lie in $[1.4422,1.5028]$), and shows that the minimum is $n-1$, attained by exactly two shapes for $n\ge 4$.
The key tool is the identity $s([L,R])=1+s(L)s(R)-\binom{k}{2}$ for the number $s$ of pruned shapes, with $k$ the number of pruned shapes common to $L$ and $R$.
\emph{Formalization:} two independent Lean proofs of the maximum, of uniqueness for $n\ge 9$ and of the minimum value.

\section*{002. The reversed-prime iteration of OEIS A100475 never enters a cycle}
\emph{Discipline: elementary number theory.}

OEIS A100475 iterates $k\mapsto\operatorname{rev}(p_k)$, where $p_k$ is the $k$-th prime and $\operatorname{rev}$ reverses decimal digits, and asks whether a start other than $1$ can lead to a loop; the question is an open statement in google-deepmind/formal-conjectures.
The note shows that no start leads to a cycle: the number of digits never decreases along an orbit, $\pi(10^d)<10^{d-1}$ for $d\ge 5$ confines any cycle to $[1,1229]$, and a finite check excludes cycles there.
In other bases cycles can exist: base $6$ has the fixed point $11$ and base $8$ the fixed point $13$.
\emph{Formalization:} a Lean proof of the negation of the formal-conjectures statement.


\section*{003. Maximally probable labeled topologies with simultaneity}
\emph{Discipline: mathematical biology (phylogenetics); combinatorics of posets.}

The tree $U^*_n$ of Harding, Hammersley and Grimmett uniquely maximizes the number of labeled histories (orderings of the internal vertices compatible with ancestry) among binary trees with $n$ leaves.
Dickey and Rosenberg (Discrete Appl.\ Math.\ 2026) allowed several branching events to occur simultaneously and conjectured that $U^*_n$ also uniquely maximizes the total number $N(T)$ of such tie-permitting histories, and maximizes the number $E(T,z)$ with $z$ events for every $z$.
The manuscript proves the conjecture in a stronger, coefficientwise form: for every bound $m$, the generating polynomial $\sum_f\prod_v x_{f(v)}$ of strict labelings of the internal vertices by $\{1,\dots,m\}$ is maximized coefficientwise by $U^*_n$.
The key tool is a switching injection comparing labelings of two trees with labelings of their union and intersection, with different label bounds on the two sides, used in an induction along the dyadic structure of $U^*_n$.
\emph{Formalization:} both parts of the conjecture including uniqueness, with the paper's recursive definition of $U^*_n$.

\end{document}
